\chapter{The Trajectory Distribution View}\label{ch:trajectory}

Chapter~\ref{ch:saliency} turns internal signals and backward routing choices into attribution maps; Chapter~\ref{ch:experiments} evaluates them. Their use in interventions also requires examining the underlying calculation. The equilibrium policies of Chapter~\ref{ch:game_theory} induce backward-trajectory distributions, enabling comparison before projection onto input pixels. We study their Hellinger distance, its exact per-input-state disagreement map, and complete versus survival-conditioned laws. Spectral analysis diagnoses amplification and target discrimination through the attribution redistribution operator's singular spectrum at a fixed input~\parencite{Maess2026Spectral}; our experiments compare normalized routing under input noise and parameter randomization.

\section{Comparison with Normalized Relevance Measures}\label{sec:nrm_comparison}

The \ac{nrm} framework of \textcite{Xiong2026nrm}, by Ping Xiong, Thomas Schnake, Grégoire Montavon, Klaus-Robert Müller, and Shinichi Nakajima, studies backward calculations through a normalized signed measure on complete neuron walks. Summing walk relevance through selected neurons or concepts gives their relevance; conditional relevance divides joint by nonzero marginal relevance. With suitable normalized propagation rules, this covers ordinary \ac{lrp}, concept relevance, and Concept Relevance Propagation (CRP)~\parencite{Achtibat2023CRP}. Their CRP quantity is joint relevance of an input neuron, intermediate concept neuron, and target output, without an automatic conditional-probability interpretation.

Both views retain network paths before forming attribution maps; their factorization supports sum-product computation and queries across layers. In \ac{nrm}, individual walk relevance may be negative despite unit total. Our game policies define a non-negative probability law $\Pi$ on the expanded graph, including player labels and termination. Signed differences of its discounted occupation measures recover gradients and relevance scores (Chapter~\ref{ch:game_theory}), generally distinct from unweighted marginals. Player expansion thus separates trajectory probability from contribution sign and scale, without identifying $\Pi$ with a given \ac{nrm} measure.

The \ac{nrm} construction normalizes propagation columns under positive-sum assumptions; identifying it with globally normalized original walk relevance additionally requires the constant-sum conditions of \textcite[Section~3.5]{Xiong2026nrm}. Our transitions remain normalized; discounts retain scale, with explicit bias and stabilizer absorption. Hellinger distance applies directly to these probability laws, but generally not to signed relevance measures. Conditioning $\Pi$ on a concept requires positive probability; signed relevance can instead vanish by cancellation.

Symbolic XAI scores logical feature queries by selecting signed terms from a multi-order prediction decomposition~\parencite{Schnake2025SymbolicXAI}. SemanticLens combines component attributions and semantic embeddings to audit expected and spurious visual concepts~\parencite{Dreyer2025SemanticLens}. These address feature relationships and human interpretation; trajectory probabilities describe normalized routing. We propose querying game trajectories through concepts, retaining attribution signs and discounts, and comparing their laws across inputs or models. Equivalence with a particular \ac{nrm}/CRP construction remains unproved and would require matching propagation rules, normalization, and absorption conventions.

\section{Trajectory Sanity Check}\label{ptnb:sec:hellinger-sanity}\label{ptnb:sec:hellinger}

The cascading-randomization sanity check of~\textcite{Adebayo2018sanity} tests explanations' dependence on learned parameters by comparing pixel-space maps. Here we examine normalized backward routing: a trajectory distribution $\Pi$ whose model-dependent signed discounted occupation projection recovers attribution.
The \textbf{Hellinger distance}~\parencite{Hellinger1909} $\mathrm{H}(\Pi_A, \Pi_B) = (1 -
    \mathrm{BC})^{1/2}$, with Bhattacharyya coefficient~\parencite{Bhattacharyya1943} $\mathrm{BC} =
    \sum_\zeta \sqrt{\Pi_A(\zeta)\,\Pi_B(\zeta)}$, directly compares routing. We test whether parameter randomization moves $\Pi$ even when attribution maps remain stable. The metric satisfies $\mathrm{H}=0$ iff $\Pi_A = \Pi_B$ and $\mathrm{H}=1$ iff their supports are disjoint. It remains finite when routes open or close, unlike KL-type divergences. For layered Markov laws on the same graph, BC factorizes: one backward pass computes $\mathrm{H}$ exactly in edge-linear time, and $\mathrm{H}^2$ decomposes into per-input-state contributions (Section~\ref{ptnb:app:path-divergence}).

\subsection{Trajectories change, maps need not}\label{ptnb:sec:hellinger-cascading}

\begin{figure}[t]
    \centering
    \begin{tikzpicture}
\begin{groupplot}[
    group style={group size=3 by 1, horizontal sep=0.40cm},
    scale only axis,
    height=2.5cm,
    xmin=1,
    ymin=0.7,
    ymax=1,
    enlarge x limits=false,
    enlarge y limits=false,
    axis x line*=bottom,
    axis y line*=left,
    axis line style={draw=black, line width=0.5pt},
    tick style={black, line width=0.45pt},
    title style={font=\footnotesize},
    ytick={0.7,0.8,0.9,1.0},
    x tick label style={font=\scriptsize},
    y tick label style={font=\scriptsize},
    tick align=outside,
    scaled x ticks=false,
    scaled y ticks=false,
    grid=both,
    major grid style={draw=gray!25},
    minor grid style={draw=gray!15},
    minor y tick num=1,
]
\nextgroupplot[title={VGG-16}, width=2.80cm, xmax=16, xtick={1,4,8,12,16}, minor xtick={1,2,3,5,6,7,9,10,11,13,14,15}]
\addplot+[color=blue!70!black, mark=none, line width=0.9pt] coordinates {(0, 0.0000000000) (1, 0.7594861872) (2, 0.9025579483) (3, 0.9087934005) (4, 0.9089968695) (5, 0.9085077148) (6, 0.9085178013) (7, 0.9079907648) (8, 0.9075277945) (9, 0.9074413281) (10, 0.9069640748) (11, 0.9069568005) (12, 0.9068374789) (13, 0.9070315441) (14, 0.9064698889) (15, 0.9064328548) (16, 0.9059700120)};
\addplot+[color=blue!70!black, mark=none, line width=0.9pt, densely dotted] coordinates {(0, 0.0000000000) (1, 0.7578458737) (2, 0.9713578362) (3, 0.9959455417) (4, 0.9998549258) (5, 0.9999820838) (6, 0.9999976776) (7, 0.9999999043) (8, 0.9999999923) (9, 0.9999999954) (10, 0.9999999954) (11, 0.9999999963) (12, 0.9999999945) (13, 0.9999999954) (14, 0.9999999953) (15, 0.9999999960) (16, 0.9999999960)};
\addplot+[color=red!80!black, mark=none, line width=0.9pt] coordinates {(0, 0.0000000000) (1, 0.7750963679) (2, 0.9771973359) (3, 0.9972827009) (4, 0.9999178271) (5, 0.9999918573) (6, 0.9999991492) (7, 0.9999999560) (8, 0.9999999848) (9, 0.9999999844) (10, 0.9999999864) (11, 0.9999999864) (12, 0.9999999857) (13, 0.9999999860) (14, 0.9999999868) (15, 0.9999999862) (16, 0.9999999852)};
\addplot+[color=red!80!black, mark=none, line width=0.9pt, densely dotted] coordinates {(0, 0.0000000000) (1, 0.7750963687) (2, 0.9771973342) (3, 0.9972826976) (4, 0.9999178316) (5, 0.9999918619) (6, 0.9999991561) (7, 0.9999999598) (8, 0.9999999909) (9, 0.9999999900) (10, 0.9999999908) (11, 0.9999999912) (12, 0.9999999913) (13, 0.9999999920) (14, 0.9999999909) (15, 0.9999999912) (16, 0.9999999937)};
\addplot+[color=red!80!black, mark=none, line width=0.45pt, solid, forget plot] coordinates {(0, 1.0000) (16, 1.0000)};
\addplot+[color=red!80!black, mark=none, line width=0.45pt, densely dotted, forget plot] coordinates {(0, 1.0000) (16, 1.0000)};
\addplot+[color=blue!70!black, mark=none, line width=0.45pt, solid, forget plot] coordinates {(0, 0.9392) (16, 0.9392)};
\addplot+[color=blue!70!black, mark=none, line width=0.45pt, densely dotted, forget plot] coordinates {(0, 1.0000) (16, 1.0000)};
\nextgroupplot[title={ResNet-50}, width=4.80cm, xmax=54, xtick={1,6,12,18,24,30,36,42,48,54}, minor xtick={1,2,3,4,5,7,8,9,10,11,13,14,15,16,17,19,20,21,22,23,25,26,27,28,29,31,32,33,34,35,37,38,39,40,41,43,44,45,46,47,49,50,51,52,53}, ylabel={}, yticklabel style={opacity=0}, ytick style={draw=none}]
\addplot+[color=blue!70!black, mark=none, line width=0.9pt] coordinates {(0, 0.0000000000) (1, 0.7740136613) (2, 0.8611538021) (3, 0.8759621524) (4, 0.8821577935) (5, 0.8949340067) (6, 0.8952563442) (7, 0.8961484951) (8, 0.8979499682) (9, 0.8996687432) (10, 0.8998905003) (11, 0.9000403845) (12, 0.8999737336) (13, 0.9004159612) (14, 0.9006932433) (15, 0.9006786146) (16, 0.9005506424) (17, 0.9005396870) (18, 0.9006573955) (19, 0.9003068308) (20, 0.9004511537) (21, 0.9009161669) (22, 0.9007179197) (23, 0.9008063574) (24, 0.9008720157) (25, 0.9010617791) (26, 0.9010646333) (27, 0.9009167788) (28, 0.9010399837) (29, 0.9010476749) (30, 0.9009851986) (31, 0.9010130783) (32, 0.9013427772) (33, 0.9013757489) (34, 0.9013422249) (35, 0.9013275622) (36, 0.9013696667) (37, 0.9013835358) (38, 0.9013705524) (39, 0.9013678672) (40, 0.9011436960) (41, 0.9011754702) (42, 0.9012944752) (43, 0.9012309719) (44, 0.9011644558) (45, 0.9011142729) (46, 0.9012432501) (47, 0.9011617513) (48, 0.9011845674) (49, 0.9009077755) (50, 0.9012393394) (51, 0.9012367828) (52, 0.9012307767) (53, 0.9012420354) (54, 0.9012335588)};
\addplot+[color=blue!70!black, mark=none, line width=0.9pt, densely dotted] coordinates {(0, 0.0000000000) (1, 0.8017318050) (2, 0.9052251460) (3, 0.9204264669) (4, 0.9278606995) (5, 0.9590658199) (6, 0.9633232114) (7, 0.9675243174) (8, 0.9941942378) (9, 0.9978053638) (10, 0.9979266594) (11, 0.9981171811) (12, 0.9984286263) (13, 0.9985279566) (14, 0.9986218022) (15, 0.9988654979) (16, 0.9989230656) (17, 0.9990260844) (18, 0.9992769944) (19, 0.9993233229) (20, 0.9994133009) (21, 0.9995320836) (22, 0.9995596989) (23, 0.9996019326) (24, 0.9996966179) (25, 0.9997284206) (26, 0.9997427923) (27, 0.9999168217) (28, 0.9999610881) (29, 0.9999647956) (30, 0.9999683051) (31, 0.9999758229) (32, 0.9999805300) (33, 0.9999808443) (34, 0.9999845910) (35, 0.9999860506) (36, 0.9999870032) (37, 0.9999901142) (38, 0.9999904376) (39, 0.9999911678) (40, 0.9999970531) (41, 0.9999979936) (42, 0.9999983322) (43, 0.9999984756) (44, 0.9999988854) (45, 0.9999989214) (46, 0.9999989972) (47, 0.9999992452) (48, 0.9999993248) (49, 0.9999994627) (50, 0.9999998364) (51, 0.9999998923) (52, 0.9999999128) (53, 0.9999999227) (54, 0.9999999902)};
\addplot+[color=red!80!black, mark=none, line width=0.9pt] coordinates {(0, 0.0000000000) (1, 0.7811890380) (2, 0.9489148586) (3, 0.9585202962) (4, 0.9586654908) (5, 0.9636757324) (6, 0.9658416452) (7, 0.9663512918) (8, 0.9643473265) (9, 0.9632637219) (10, 0.9633963239) (11, 0.9633992485) (12, 0.9665420467) (13, 0.9672464530) (14, 0.9676737033) (15, 0.9688082100) (16, 0.9687992955) (17, 0.9691870298) (18, 0.9701412473) (19, 0.9703885827) (20, 0.9708043630) (21, 0.9715813133) (22, 0.9715691819) (23, 0.9718293913) (24, 0.9721858271) (25, 0.9725370294) (26, 0.9726692293) (27, 0.9720056313) (28, 0.9725633618) (29, 0.9726720279) (30, 0.9726074320) (31, 0.9740267482) (32, 0.9749823267) (33, 0.9750237942) (34, 0.9755110628) (35, 0.9756800647) (36, 0.9757804705) (37, 0.9762739732) (38, 0.9764651931) (39, 0.9764496777) (40, 0.9757507141) (41, 0.9757533700) (42, 0.9762824471) (43, 0.9757907308) (44, 0.9771094255) (45, 0.9768832311) (46, 0.9771116252) (47, 0.9771881909) (48, 0.9773463414) (49, 0.9772364447) (50, 0.9778191962) (51, 0.9778556043) (52, 0.9778459094) (53, 0.9776280403) (54, 0.9782296548)};
\addplot+[color=red!80!black, mark=none, line width=0.9pt, densely dotted] coordinates {(0, 0.0000000000) (1, 0.7829772820) (2, 0.9798537486) (3, 0.9895289502) (4, 0.9909433136) (5, 0.9959472404) (6, 0.9973671432) (7, 0.9976942099) (8, 0.9990083083) (9, 0.9997192996) (10, 0.9998309772) (11, 0.9998639833) (12, 0.9999371075) (13, 0.9999581831) (14, 0.9999650429) (15, 0.9999821148) (16, 0.9999872373) (17, 0.9999895926) (18, 0.9999948138) (19, 0.9999963344) (20, 0.9999971676) (21, 0.9999985077) (22, 0.9999988877) (23, 0.9999990597) (24, 0.9999994747) (25, 0.9999996324) (26, 0.9999996750) (27, 0.9999998401) (28, 0.9999999490) (29, 0.9999999640) (30, 0.9999999719) (31, 0.9999999828) (32, 0.9999999878) (33, 0.9999999864) (34, 0.9999999877) (35, 0.9999999896) (36, 0.9999999891) (37, 0.9999999901) (38, 0.9999999900) (39, 0.9999999895) (40, 0.9999999887) (41, 0.9999999882) (42, 0.9999999912) (43, 0.9999999899) (44, 0.9999999902) (45, 0.9999999919) (46, 0.9999999905) (47, 0.9999999913) (48, 0.9999999900) (49, 0.9999999887) (50, 0.9999999900) (51, 0.9999999899) (52, 0.9999999903) (53, 0.9999999893) (54, 0.9999999901)};
\addplot+[color=red!80!black, mark=none, line width=0.45pt, solid, forget plot] coordinates {(0, 0.9728) (54, 0.9728)};
\addplot+[color=red!80!black, mark=none, line width=0.45pt, densely dotted, forget plot] coordinates {(0, 1.0000) (54, 1.0000)};
\addplot+[color=blue!70!black, mark=none, line width=0.45pt, solid, forget plot] coordinates {(0, 0.9382) (54, 0.9382)};
\addplot+[color=blue!70!black, mark=none, line width=0.45pt, densely dotted, forget plot] coordinates {(0, 1.0000) (54, 1.0000)};
\nextgroupplot[title={ViT-B/16}, width=3.70cm, xmax=38, xtick={1,6,12,18,24,30,36}, minor xtick={1,2,3,4,5,7,8,9,10,11,13,14,15,16,17,19,20,21,22,23,25,26,27,28,29,31,32,33,34,35,37}, ylabel={}, yticklabel style={opacity=0}, ytick style={draw=none}, legend style={font=\scriptsize, draw=black, fill=white, fill opacity=0.92, text opacity=1, cells={anchor=west}, row sep=1pt, inner xsep=4pt, inner ysep=3pt, at={(0.98,0.04)}, anchor=south east}, legend columns=1]
\addplot+[color=blue!70!black, mark=none, line width=0.9pt] coordinates {(0, 0.0000000000) (1, 0.7532651744) (2, 0.8319595623) (3, 0.8478672457) (4, 0.8844686497) (5, 0.9110840737) (6, 0.9180494944) (7, 0.9297589499) (8, 0.9392713147) (9, 0.9420341653) (10, 0.9459173348) (11, 0.9487002171) (12, 0.9495841725) (13, 0.9507822544) (14, 0.9518463139) (15, 0.9520969598) (16, 0.9526289120) (17, 0.9529080283) (18, 0.9526933004) (19, 0.9529460436) (20, 0.9531129656) (21, 0.9531364352) (22, 0.9532026722) (23, 0.9530297901) (24, 0.9522956489) (25, 0.9523376464) (26, 0.9522885949) (27, 0.9518661545) (28, 0.9518677129) (29, 0.9519770299) (30, 0.9513401467) (31, 0.9513218235) (32, 0.9512989790) (33, 0.9503164608) (34, 0.9503572194) (35, 0.9505905549) (36, 0.9501477794) (37, 0.9503059207) (38, 0.9502751929)};
\addlegendentry{SG $\mathrm{H}$}
\addplot+[color=blue!70!black, mark=none, line width=0.9pt, densely dotted] coordinates {(0, 0.0000000000) (1, 0.7534678910) (2, 0.7582185272) (3, 0.8041403636) (4, 0.8926397729) (5, 0.9016008099) (6, 0.9204114223) (7, 0.9550356159) (8, 0.9605720967) (9, 0.9685903810) (10, 0.9823387799) (11, 0.9833015107) (12, 0.9866931981) (13, 0.9923258868) (14, 0.9930370014) (15, 0.9944731765) (16, 0.9968666783) (17, 0.9971581010) (18, 0.9976668187) (19, 0.9986798657) (20, 0.9988496761) (21, 0.9990755547) (22, 0.9994669202) (23, 0.9995103887) (24, 0.9995887385) (25, 0.9997593445) (26, 0.9997910467) (27, 0.9998358110) (28, 0.9999014361) (29, 0.9999208442) (30, 0.9999344799) (31, 0.9999592826) (32, 0.9999663043) (33, 0.9999718485) (34, 0.9999812348) (35, 0.9999860119) (36, 0.9999875562) (37, 0.9999911246) (38, 0.9999965758)};
\addlegendentry{SG $\mathrm{H}_{\mathrm{surv}}$}
\addplot+[color=red!80!black, mark=none, line width=0.9pt] coordinates {(0, 0.0000000000) (1, 0.7679664488) (2, 0.8507813274) (3, 0.8843068618) (4, 0.9149508559) (5, 0.9402741843) (6, 0.9493400595) (7, 0.9576074540) (8, 0.9651670280) (9, 0.9680426625) (10, 0.9704187923) (11, 0.9723567481) (12, 0.9727836145) (13, 0.9734118713) (14, 0.9740074796) (15, 0.9735864876) (16, 0.9737970057) (17, 0.9739868146) (18, 0.9729427524) (19, 0.9731612574) (20, 0.9733292292) (21, 0.9722940493) (22, 0.9723583706) (23, 0.9724112122) (24, 0.9702854529) (25, 0.9703439373) (26, 0.9703494959) (27, 0.9680573755) (28, 0.9681142773) (29, 0.9681452230) (30, 0.9659158272) (31, 0.9660585551) (32, 0.9661357114) (33, 0.9628330510) (34, 0.9628127877) (35, 0.9636050890) (36, 0.9618825492) (37, 0.9622859025) (38, 0.9626814679)};
\addlegendentry{RG $\mathrm{H}$}
\addplot+[color=red!80!black, mark=none, line width=0.9pt, densely dotted] coordinates {(0, 0.0000000000) (1, 0.7685156618) (2, 0.7981041034) (3, 0.8566331953) (4, 0.9033373994) (5, 0.9208234133) (6, 0.9434315864) (7, 0.9611812232) (8, 0.9684138554) (9, 0.9782742332) (10, 0.9853145354) (11, 0.9872186758) (12, 0.9914135539) (13, 0.9942586743) (14, 0.9951891577) (15, 0.9966983554) (16, 0.9978136524) (17, 0.9981886543) (18, 0.9987189387) (19, 0.9991583191) (20, 0.9993203225) (21, 0.9995171486) (22, 0.9996765463) (23, 0.9997372367) (24, 0.9998044894) (25, 0.9998696226) (26, 0.9998993841) (27, 0.9999248884) (28, 0.9999494906) (29, 0.9999641580) (30, 0.9999721711) (31, 0.9999809615) (32, 0.9999861521) (33, 0.9999892142) (34, 0.9999925991) (35, 0.9999951875) (36, 0.9999960188) (37, 0.9999973923) (38, 0.9999992663)};
\addlegendentry{RG $\mathrm{H}_{\mathrm{surv}}$}
\addplot+[color=red!80!black, mark=none, line width=0.45pt, solid, forget plot] coordinates {(0, 0.9563) (38, 0.9563)};
\addplot+[color=red!80!black, mark=none, line width=0.45pt, densely dotted, forget plot] coordinates {(0, 1.0000) (38, 1.0000)};
\addplot+[color=blue!70!black, mark=none, line width=0.45pt, solid, forget plot] coordinates {(0, 0.9579) (38, 0.9579)};
\addplot+[color=blue!70!black, mark=none, line width=0.45pt, densely dotted, forget plot] coordinates {(0, 1.0000) (38, 1.0000)};
\end{groupplot}
\end{tikzpicture}

    \caption[Hellinger trajectory distance under parameter randomization]{Hellinger trajectory distance under cascading parameter
    randomization~\parencite{Adebayo2018sanity}. Hellinger distance $\mathrm{H}$ (vertical axis, from
    $0.7$) against the number of cascadingly re-initialized layers, output to input (horizontal axis;
    $\mathrm{H}=0$ before the first layer is randomized). Solid: plain $\mathrm{H}$; dotted:
    $\mathrm{H}_{\mathrm{surv}}$, conditioned on survival to an ordinary input-layer state. Thin
    horizontal reference lines mark the double-random baseline values from
    Table~\ref{ptnb:tab:double-random-hellinger} (solid: plain $\mathrm{H}$; dotted:
    $\mathrm{H}_{\mathrm{surv}}$), in the color of the corresponding game.}
    \label{ptnb:fig:hellinger-sanity}
\end{figure}

Under cascading randomization (layers cumulatively re-initialized output$\to$input), \ac{sg} and \ac{rg} diverge from the first randomized layer and plateau near the double-random reference (Table~\ref{ptnb:tab:double-random-hellinger}, Figure~\ref{ptnb:fig:hellinger-sanity}). This reference lies below $\mathrm{H}=1$ because models randomized over the tested layers still share short trajectories: both may stop at structurally dead neurons and enter the common cemetery. Such paths coincide more readily than full input-reaching trajectories. Conditioning on survival removes shared cemetery mass, yielding $\mathrm{H}_{\mathrm{surv}} \approx 1$. Trajectory and map readouts have weak observed per-image rank association (Spearman $|\rho|\le \perImageCrossMax$ between $\mathrm{H}$ and pixel-level Gradient or $\alpha\beta$-\ac{lrp} similarity): similar maps need not imply similar calculations. Section~\ref{ptnb:app:sanity-check} details per-game values, conditioning, the floor mechanism, and per-image analysis.

\paragraph{What the map hides.}\label{ptnb:sec:pixel-sanity}
Map-level checks under parameter randomization reproduce~\textcite{Adebayo2018sanity} and \textcite{sixt2020whenexplanationslie}: some methods retain structure at the cascade endpoint. On ViT-B/16, the \ac{rg} anchor retains Spearman $|\rho| = \sanityVitRGanchorRho$, highest among all methods, and HOG $r = \sanityVitRGanchorHOG$; VGG-16 LRP-$\varepsilon$ retains $\sanityVggLRPepsRho$ and $\sanityVggLRPepsHOG$ (Table~\ref{tab:sanity-check}). Localization-selected \ac{rg} modes are slightly more model sensitive than the anchor on all three models. \ch{In Figure~\ref{ptnb:fig:hellinger-sanity}, $\mathrm{H}$ reaches the double-random reference while maps retain structure. The map-level failure remains; $\mathrm{H}$ detects routing changes obscured by the map. Since discounted occupation projections also vary across models, this discrepancy does not isolate loss under a fixed projection.}
\section{Hellinger Distance and Disagreement Map}\label{ptnb:app:path-divergence}

We compare two layered Markov processes on a common state graph and seek both a scalar divergence
of the trajectory distributions and an exact per-pixel decomposition of this divergence at the
input layer.  We use the squared Hellinger distance because the geometric mean $\sqrt{PQ}$
preserves the multiplicative structure of Markov trajectory distributions and yields an exact
linear-time dynamic program.  Unlike Jeffreys divergence, it remains finite under support mismatch;
unlike Jensen--Shannon, it admits an exact recursion on the layered graph.  Later, we decompose the
squared Hellinger distance as a sum of non-negative per-pixel contributions over all input-layer
states, giving a disagreement map at the input layer (Theorem~\ref{ptnb:thm:per-pixel}).

The formal foundations and proofs are given in companion Chapter~\ref{app:supp_trajectory}.
We compare two processes on a common finite layered state graph with a shared deterministic
starting state $s_0$. At layer $l$, $\mathcal S^{(l)}_+=\mathcal S^{(l)}\cup\{\dagger\}$
contains the ordinary states and an absorbing cemetery state; $T_M^{(l)}$ is the transition
kernel from layer $l$ to $l-1$, and $m_M^{(l)}(s)$ is the marginal probability of visiting $s$.
The full trajectory law is $\Pi_M=\Pi_M^{(0)}$.
For the \ac{rg}, full-law diagnostics retain the gate, branch, and routing substeps; only
ordinary surviving paths admit exact act-only elimination
(Remark~\ref{ptnb:rem:rg-natural-equivalence}).

\subsection{Bhattacharyya Coefficient and the Geometric-Mean Kernel}\label{ptnb:sec:lmp-bc}

\paragraph{Definitions.}
For two probability distributions $P, Q$ on a common finite set, the \emph{Bhattacharyya
coefficient} is the geometric overlap
\begin{align}\label{ptnb:eq:bc-def}
    \mathrm{BC}(P, Q) \;\coloneqq\; \sum_{y} \sqrt{P(y)\, Q(y)} \;\in\; [0, 1],
\end{align}
and the \emph{squared Hellinger distance} is
\begin{align}\label{ptnb:eq:hellinger-def}
    \mathrm{H}^2(P, Q) \;\coloneqq\; \tfrac{1}{2} \sum_y \bigl(\sqrt{P(y)} - \sqrt{Q(y)}\bigr)^2 \;=\;
    1 - \mathrm{BC}(P, Q) \;\in\; [0, 1].
\end{align}
The Hellinger distance $\mathrm{H}(P,Q) = \sqrt{1 - \mathrm{BC}(P,Q)}$ is a true metric: $0 \leq
    \mathrm{H} \leq 1$, with $\mathrm{H} = 0$ iff $P = Q$ and $\mathrm{H} = 1$ iff $P$ and $Q$ have
disjoint supports.  Crucially, $\sqrt{0 \cdot Q(y)} = 0$, so mismatched supports cause no
infinity---only a smaller overlap.

\paragraph{Geometric-mean kernel.}
For each layer $l \geq 1$ and pair of states $(s, s') \in \mathcal{S}^{(l)}_+ \times
    \mathcal{S}^{(l-1)}_+$, define the \emph{geometric-mean kernel}
\begin{align}\label{ptnb:eq:geom-kernel}
    \widetilde{T}^{(l)}(s, s') \;\coloneqq\; \sqrt{T^{(l)}_A(s \to s')\; T^{(l)}_B(s \to s')}.
\end{align}
This kernel is non-negative but in general sub-stochastic.  Its row sum
\begin{align}\label{ptnb:eq:geom-rowsum}
    \widetilde{R}^{(l)}(s) \;\coloneqq\; \sum_{s' \in \mathcal{S}^{(l-1)}_+} \widetilde{T}^{(l)}(s, s')
\end{align}
satisfies $0 \leq \widetilde{R}^{(l)}(s) \leq 1$ by Cauchy--Schwarz, with equality iff $T^{(l)}_A(s
    \to \cdot) = T^{(l)}_B(s \to \cdot)$.  Cemetery handling is automatic: $T^{(l)}_A(\dagger \to
    \dagger) = T^{(l)}_B(\dagger \to \dagger) = 1$, so $\widetilde{T}^{(l)}(\dagger, \dagger) = 1$ and
$\widetilde{R}^{(l)}(\dagger) = 1$.

\subsection{Hellinger backward pass}\label{ptnb:sec:lmp-hellinger-backward}

The chain rule (Lemma~\ref{ptnb:lem:bc-chain-rule}) yields a backward DP that propagates the per-state
cumulative $\widetilde{T}$-mass $g^{(l)}(s)$ from the start state $s_0$ down through the layers.
The prefix Bhattacharyya coefficient at each layer is the row sum $\mathrm{BC}^{(l)} = \sum_{s}
    g^{(l)}(s)$, the prefix Hellinger distance is $\mathrm{H}^{(l)} = \sqrt{1 - \mathrm{BC}^{(l)}}$,
and the per-terminal-state disagreement map is read off at $l = 0$ from~\eqref{ptnb:eq:per-pixel-h2}.
The map is defined on the full terminal set $\mathcal{S}^{(0)}_+ = \mathcal{S}^{(0)} \cup
    \{\dagger\}$; the visualized pixel attribution map is its restriction to ordinary pixels
$\mathcal{S}^{(0)}$, and its sum is therefore $\mathrm{H}^2_{\mathrm{ord}} = \mathrm{H}^2 -
    \mathrm{H}^2(\dagger)$, not the full $\mathrm{H}^2$ (Theorem~\ref{ptnb:thm:per-pixel}(iii)).

\begin{algorithm}[H]
    \caption{Hellinger backward pass (with per-terminal-state disagreement map)}
    \label{ptnb:alg:hellinger}\label{ptnb:alg:hellinger-backward}
    Initialize the occupation marginals $m_M^{(L)}$ and overlap mass $g^{(L)}$ at the target output.\;
For $l=L,L-1,\ldots,1$, propagate each model's marginals with its kernel and propagate overlap by
$g^{(l-1)}(s')=\sum_s g^{(l)}(s)\sqrt{T_A^{(l)}(s,s')T_B^{(l)}(s,s')}$.\;
At each layer compute $\mathrm{BC}^{(l)}=\sum_s g^{(l)}(s)$ and
$\mathrm{H}^{(l)}=\sqrt{1-\mathrm{BC}^{(l)}}$.\;
At every terminal state return
$\mathrm{H}^2(s)=\tfrac12(m_A^{(0)}(s)+m_B^{(0)}(s))-g^{(0)}(s)$.
The pixel disagreement map restricts this output to ordinary input states.

\end{algorithm}

\subsection{Per-Pixel Disagreement Map}\label{ptnb:sec:lmp-disagreement}

The squared Hellinger distance $\mathrm{H}^2$ admits an exact non-negative decomposition over
input-layer states: each input pixel carries a calibrated share of the total trajectory
disagreement (Theorem~\ref{ptnb:thm:per-pixel}), giving a disagreement \emph{map} that is itself an
attribution of the trajectory mismatch back to the input. The conditioned-on-survival variant
removes the cemetery contribution that otherwise bounds plain $\mathrm{H}$ from below on randomized
models.

\subsubsection{Per-Pixel Decomposition: The Disagreement Map}\label{ptnb:sec:lmp-per-pixel}

The squared Hellinger distance at layer~$0$ decomposes exactly over terminal pixels.

\begin{theoremenv}[Per-pixel decomposition, \LeanResult{ptnb:thm:per-pixel}]\label{ptnb:thm:per-pixel}
Let $\mathrm{H}^2 \coloneqq \mathrm{H}^2\!\bigl(\Pi^{(0)}_A,\, \Pi^{(0)}_B\bigr)$.  For $s \in
    \mathcal{S}^{(0)}_+$, define
\begin{align}\label{ptnb:eq:per-pixel-h2}
    \mathrm{H}^2(s) \;\coloneqq\; \tfrac{1}{2}\bigl(m^{(0)}_A(s) + m^{(0)}_B(s)\bigr) \;-\; g^{(0)}(s),
\end{align}
where $g^{(0)}(s)$ is the cumulative $\widetilde{T}$-mass at $s$.  Then:
\begin{enumerate}
    \item[(i)] $\mathrm{H}^2(s) \geq 0$ for every $s \in \mathcal{S}^{(0)}_+$.
    \item[(ii)]
        \begin{align}\label{ptnb:eq:per-pixel-sum}
            \mathrm{H}^2 \;=\; \sum_{s \in \mathcal{S}^{(0)}_+} \mathrm{H}^2(s).
        \end{align}
    \item[(iii)]
        \begin{align}\label{ptnb:eq:ordinary-h2}
            \mathrm{H}^2_{\mathrm{ord}} \;\coloneqq\; \sum_{s \in \mathcal{S}^{(0)}} \mathrm{H}^2(s) \;=\;
            \mathrm{H}^2 - \mathrm{H}^2(\dagger)
        \end{align}
\end{enumerate}\end{theoremenv}

Proofs of the backward pass and the decomposition are given in companion Chapter~\ref{app:supp_trajectory}.

\paragraph{The disagreement map.}
Theorem~\ref{ptnb:thm:per-pixel} defines $\mathrm{H}^2(\cdot)$ on the full terminal set
$\mathcal{S}^{(0)}_+ = \mathcal{S}^{(0)} \cup \{\dagger\}$: this is the \emph{full terminal-state
decomposition}, with $\sum_{s \in \mathcal{S}^{(0)}_+} \mathrm{H}^2(s) = \mathrm{H}^2_{(0)}$, and
the cemetery contribution $\mathrm{H}^2(\dagger)$ is part of it. The \emph{Hellinger disagreement
map} is the ordinary-pixel restriction $s \mapsto \mathrm{H}^2(s)$ on $\mathcal{S}^{(0)}$, which is
the quantity we visualize. It is non-negative, sums to $\mathrm{H}^2_{\mathrm{ord}}$ over ordinary
pixels, and is computed from $g^{(0)}, m_A^{(0)}, m_B^{(0)}$.  The quantity
$\mathrm{H}^2_{\mathrm{ord}}$ is an unnormalized ordinary-state subsum; the conditioned-on-survival
Hellinger distance that normalizes it is introduced below. Recall from Section~\ref{ptnb:sec:games-fit}
that layer-$0$ ordinary states of the two games are \emph{tagged} pixel pairs $s^{(0)}_{j,\pm}$,
tagged by the player $q$ for the \ac{sg} and by the stream sign $\kappa$ for the \ac{rg}; the
per-pixel attribution displayed in figures aggregates the two tags,
\begin{align}
    \mathrm{H}^2_{\mathrm{pixel}}(j) \;=\; \mathrm{H}^2\bigl(s^{(0)}_{j,+}\bigr) +
    \mathrm{H}^2\bigl(s^{(0)}_{j,-}\bigr),
\end{align}
which remains non-negative and sums to $\mathrm{H}^2_{\mathrm{ord}}$ over pixels.  By the
Cauchy--Schwarz step in the proof of Theorem~\ref{ptnb:thm:per-pixel}(i), $g^{(0)}(s) \leq
    \sqrt{m^{(0)}_A(s)\,m^{(0)}_B(s)}$, so $\mathrm{H}^2(s)$ is bounded below by the marginal surrogate
$\tfrac{1}{2}\bigl(m^{(0)}_A(s) + m^{(0)}_B(s)\bigr) - \sqrt{m^{(0)}_A(s)\,m^{(0)}_B(s)}$, the
per-pixel squared Hellinger distance between the layer-0 marginal occupations; the trajectory map
is strictly larger whenever the conditional distributions of trajectories reaching pixel~$s$ differ
between models.

\subsubsection{Conditioning on Survival to the Ordinary Input
Layer}\label{ptnb:sec:lmp-conditioned-survival}

When the cemetery absorbs a large fraction of trajectory mass, the total Hellinger distance mixes
two effects: disagreement about \emph{whether} a trajectory survives and disagreement about
\emph{where surviving trajectories go}.  To isolate the second effect, we condition each model's
trajectory law on the event that the trajectory reaches an ordinary input-layer state.  The
resulting scalar compares the two \emph{surviving} trajectory distributions directly and removes
the common cemetery floor.  Because survival to the input layer factorizes along the layered graph,
this conditioning remains Markov and is implemented by a per-state Doob $h$-transform (with $h =
    \psi_M$, the remaining survival probability).

Let $\mathsf{Surv}$ denote this event and assume that both survival masses are positive. With
\[
 Z_M=\mathbb P_M(\mathsf{Surv})
     =\sum_{s\in\mathcal S^{(0)}}m_M^{(0)}(s)>0,
 \qquad
 \mathrm{BC}_{\mathrm{ord}}^{(0)}=\sum_{s\in\mathcal S^{(0)}}g^{(0)}(s),
\]
the conditioned-survival Hellinger distance is
\[
 \mathrm H_{\mathrm{surv}}^2
 =\mathrm H^2\!\bigl(\Pi_A(\,\cdot\mid\mathsf{Surv}),
                         \Pi_B(\,\cdot\mid\mathsf{Surv})\bigr)
 =1-\frac{\mathrm{BC}_{\mathrm{ord}}^{(0)}}{\sqrt{Z_A Z_B}}.
\]
Theorem~\ref{ptnb:thm:conditioned-survival} gives the conditioned kernels and the proof;
Equation~\eqref{ptnb:eq:conditioned-per-pixel-h2} gives the corresponding non-negative map.
Post-hoc normalization becomes numerically infeasible for deep networks when either
survival probability gets too small; the conditioned-kernel computation avoids this final division.

\subsection{Experiments}\label{ptnb:sec:lmp-experiments}

With the model fixed, we test sensitivity to input perturbations and separation among ImageNet-S classes and from random noise. We use the Stopping-Game kernel~\eqref{ptnb:eq:adf-T} without risk aversion ($\lambda = 0$; ``ADF'' in code) and Routing-Game kernel~\eqref{ptnb:eq:rg-T-act} at $\tau = 1$. The diagnostic \ac{rg} omits the bias action and renormalizes the remaining branches, so it may differ from the bias-inclusive attribution game.

\subsubsection{Input-Noise Sensitivity}\label{ptnb:sec:input-noise-hellinger}

For each architecture and $\sigma \in \{0.005,\allowbreak 0.01,\allowbreak 0.02,\allowbreak 0.03,
    \allowbreak 0.05,\allowbreak 0.1,\allowbreak 0.2,\allowbreak 0.3,\allowbreak 0.5\}$, we draw one standard-Gaussian tensor $\eta$ per image and compute $\mathrm{H}(\Pi_A(x), \Pi_A(x + \sigma\eta))$ for both kernels, targeting the clean input's argmax. Mean $\mathrm{H}$ and $\mathrm{H}_{\mathrm{surv}}$ increase with $\sigma$ for both games on all three architectures (Figure~\ref{ptnb:fig:hellinger-input-noise}). They coincide for VGG-16 \ac{rg}, whose kernel never routes to dead-neuron states. ResNet-50 and ViT-B/16 skip connections reach such states; ViT-B/16's GELU gate also sends some trajectories to the cemetery at every state. Generalised Explanation Faithfulness (GEF) correlates explanation and model-output distortion along perturbation paths~\parencite{Hedstrom2025GEF}, anchoring explanation changes to model responses. Our Hellinger curves measure routing sensitivity, without establishing this alignment or explanation faithfulness.

\begin{figure}[H]
    \centering
    \begin{tikzpicture}
\begin{groupplot}[
    group style={group size=3 by 1, horizontal sep=0.90cm},
    scale only axis,
    height=3.25cm,
    ymin=0,
    enlarge x limits=false,
    enlarge y limits=false,
    axis x line*=bottom,
    axis y line*=left,
    axis line style={draw=black, line width=0.8pt},
    tick style={black, line width=0.7pt},
    xlabel={Input-noise $\sigma$},
    ylabel={Hellinger distance $\mathrm{H}$},
    xlabel style={font=\footnotesize},
    ylabel style={font=\footnotesize},
    title style={font=\footnotesize},
    x tick label style={font=\scriptsize},
    y tick label style={font=\scriptsize},
    tick align=outside,
    scaled x ticks=false,
    scaled y ticks=false,
    grid=both,
    major grid style={draw=gray!25},
    minor grid style={draw=gray!15},
    minor y tick num=1,
]
\nextgroupplot[title={VGG-16}, width=3.60cm, xmode=log, log basis x=10, xmax=0.500, xmin=0.0045, log ticks with fixed point, xtick={0.005, 0.01, 0.02, 0.03, 0.05, 0.1, 0.2, 0.3, 0.5}, x tick label style={rotate=90, anchor=east, /pgf/number format/.cd, fixed, precision=3}, ymax=1.0, ytick={0,0.2,0.4,0.6,0.8,1.0}, yticklabel style={font=\scriptsize, /pgf/number format/.cd, fixed, precision=1, zerofill}]
\addplot+[color=blue!70!black, mark=none, line width=1.5pt] coordinates {(0.005000, 0.0553140999) (0.010000, 0.0856729081) (0.020000, 0.1308442531) (0.030000, 0.1647212310) (0.050000, 0.2165156573) (0.100000, 0.2942934630) (0.200000, 0.3756936446) (0.300000, 0.4278060399) (0.500000, 0.5033394820)};
\addplot+[color=blue!70!black, mark=none, line width=1.5pt, densely dotted] coordinates {(0.005000, 0.4797141252) (0.010000, 0.6291917364) (0.020000, 0.7834284200) (0.030000, 0.8614684689) (0.050000, 0.9347354570) (0.100000, 0.9844335426) (0.200000, 0.9978674384) (0.300000, 0.9995125510) (0.500000, 0.9999488152)};
\addplot+[color=red!80!black,  mark=none, line width=1.5pt] coordinates {(0.005000, 0.2776190381) (0.010000, 0.3922999387) (0.020000, 0.5496846290) (0.030000, 0.6581073065) (0.050000, 0.7958674986) (0.100000, 0.9337734384) (0.200000, 0.9888058215) (0.300000, 0.9973915839) (0.500000, 0.9997982972)};
\addplot+[color=red!80!black, mark=none, line width=1.5pt, densely dotted] coordinates {(0.005000, 0.2775959566) (0.010000, 0.3922755395) (0.020000, 0.5496613847) (0.030000, 0.6580863276) (0.050000, 0.7958493750) (0.100000, 0.9337572963) (0.200000, 0.9887889032) (0.300000, 0.9973742499) (0.500000, 0.9997809263)};
\nextgroupplot[title={ResNet-50}, width=3.60cm, xmode=log, log basis x=10, xmax=0.500, xmin=0.0045, log ticks with fixed point, xtick={0.005, 0.01, 0.02, 0.03, 0.05, 0.1, 0.2, 0.3, 0.5}, x tick label style={rotate=90, anchor=east, /pgf/number format/.cd, fixed, precision=3}, ymax=1.0, ytick={0,0.2,0.4,0.6,0.8,1.0}, ylabel={}, yticklabel style={font=\scriptsize, /pgf/number format/.cd, fixed, precision=1, zerofill}]
\addplot+[color=blue!70!black, mark=none, line width=1.5pt] coordinates {(0.005000, 0.0452152578) (0.010000, 0.0642118959) (0.020000, 0.0907002623) (0.030000, 0.1095971441) (0.050000, 0.1325101551) (0.100000, 0.1642478283) (0.200000, 0.1974226976) (0.300000, 0.2196046893) (0.500000, 0.2537678290)};
\addplot+[color=blue!70!black, mark=none, line width=1.5pt, densely dotted] coordinates {(0.005000, 0.3833659495) (0.010000, 0.5009903215) (0.020000, 0.6361947477) (0.030000, 0.7155634523) (0.050000, 0.8049244294) (0.100000, 0.8970243870) (0.200000, 0.9539478054) (0.300000, 0.9741317490) (0.500000, 0.9894004723)};
\addplot+[color=red!80!black,  mark=none, line width=1.5pt] coordinates {(0.005000, 0.1178935577) (0.010000, 0.1738609446) (0.020000, 0.2586063418) (0.030000, 0.3201846700) (0.050000, 0.3993576777) (0.100000, 0.4971650731) (0.200000, 0.5790294239) (0.300000, 0.6243918148) (0.500000, 0.6840088309)};
\addplot+[color=red!80!black, mark=none, line width=1.5pt, densely dotted] coordinates {(0.005000, 0.2716503111) (0.010000, 0.3841185327) (0.020000, 0.5427390559) (0.030000, 0.6517284764) (0.050000, 0.7838204504) (0.100000, 0.9166907404) (0.200000, 0.9796536898) (0.300000, 0.9933897417) (0.500000, 0.9989774765)};
\nextgroupplot[title={ViT-B/16}, width=3.60cm, xmode=log, log basis x=10, xmax=0.500, xmin=0.0045, log ticks with fixed point, xtick={0.005, 0.01, 0.02, 0.03, 0.05, 0.1, 0.2, 0.3, 0.5}, x tick label style={rotate=90, anchor=east, /pgf/number format/.cd, fixed, precision=3}, ymax=1.0, ytick={0,0.2,0.4,0.6,0.8,1.0}, ylabel={}, yticklabel style={font=\scriptsize, /pgf/number format/.cd, fixed, precision=1, zerofill}, legend style={font=\scriptsize, draw=black, fill=white, fill opacity=0.92, text opacity=1, cells={anchor=west}, row sep=1pt, inner xsep=4pt, inner ysep=3pt, at={(0.02,0.98)}, anchor=north west}, legend columns=1]
\addplot+[color=blue!70!black, mark=none, line width=1.5pt] coordinates {(0.005000, 0.0133444159) (0.010000, 0.0265228240) (0.020000, 0.0543220247) (0.030000, 0.0813419521) (0.050000, 0.1233459404) (0.100000, 0.1851983074) (0.200000, 0.2442760116) (0.300000, 0.2769362447) (0.500000, 0.3133687389)};
\addlegendentry{SG $\mathrm{H}$}
\addplot+[color=blue!70!black, mark=none, line width=1.5pt, densely dotted] coordinates {(0.005000, 0.0329339274) (0.010000, 0.0657790490) (0.020000, 0.1341236114) (0.030000, 0.2007487602) (0.050000, 0.3038745886) (0.100000, 0.4582102226) (0.200000, 0.6060417373) (0.300000, 0.6873592496) (0.500000, 0.7752400235)};
\addlegendentry{SG $\mathrm{H}_{\mathrm{surv}}$}
\addplot+[color=red!80!black,  mark=none, line width=1.5pt] coordinates {(0.005000, 0.0274186654) (0.010000, 0.0522607684) (0.020000, 0.1018637235) (0.030000, 0.1512062222) (0.050000, 0.2300948891) (0.100000, 0.3514355790) (0.200000, 0.4755279295) (0.300000, 0.5472851915) (0.500000, 0.6320482010)};
\addlegendentry{RG $\mathrm{H}$}
\addplot+[color=red!80!black, mark=none, line width=1.5pt, densely dotted] coordinates {(0.005000, 0.0637267384) (0.010000, 0.1212389912) (0.020000, 0.2323997733) (0.030000, 0.3368058769) (0.050000, 0.4977321302) (0.100000, 0.7227154583) (0.200000, 0.8925281441) (0.300000, 0.9512227121) (0.500000, 0.9858833916)};
\addlegendentry{RG $\mathrm{H}_{\mathrm{surv}}$}
\end{groupplot}
\end{tikzpicture}

    \caption[Input-noise Hellinger trajectory distance]{Input-noise Hellinger trajectory distance. Solid: plain $\mathrm{H}$; dotted:
    $\mathrm{H}_{\mathrm{surv}}$, conditioned on survival to an ordinary input-layer state. $x$-axis is
    log-scale input-noise $\sigma$; same $(A,A)$ model on both sides, only the input differs.}
    \label{ptnb:fig:hellinger-input-noise}
\end{figure}

\subsubsection{Class Separability and Reference Baselines}\label{ptnb:sec:input-noise-baselines}

Class separation depends on architecture and kernel (Table~\ref{ptnb:tab:hellinger-input-random}):
\begin{itemize}
    \item On \textbf{VGG-16}, \ac{sg} gives clean-intraclass $\mathrm{H}=0.513$ below clean-interclass $\mathrm{H}=0.607$. \ac{rg} saturates at $\mathrm{H}=1$ throughout, resolving no class structure.
    \item On \textbf{ResNet-50}, \ac{rg} separates clean-intraclass $\mathrm{H}=0.786$ from clean-interclass $\mathrm{H}=0.902$; \ac{sg} does not (intraclass $0.272$ vs.\ interclass $0.268$).
    \item On \textbf{ViT-B/16}, \ac{rg} gives intraclass $0.789$ vs.\ interclass $0.882$; the \ac{sg} gap ($0.422 \pm 0.023$ vs.\ $0.435 \pm 0.017$) is smaller than the spread over pairs and does not separate classes.
\end{itemize}
Random-vs-random is lowest in every nonconstant column. Clean-vs-random is highest in each plain-$\mathrm{H}$ column except ResNet-50 \ac{sg} ($0.269$ vs.\ intraclass $0.272$) and saturated VGG-16 \ac{rg}, consistent with different routing for Gaussian and in-distribution inputs. Conditioned $\mathrm{H}_{\mathrm{surv}}$ is at least $0.99$ throughout VGG-16, ResNet-50, and ViT-B/16 \ac{rg}, providing no class-separability signal. ViT-B/16 \ac{sg} separates random-vs-random ($0.852$) from clean rows ($0.957$--$0.968$), but neither intraclass from interclass ($0.963$ vs.\ $0.968$) nor clean-vs-random ($0.957$) from clean pairs.

\begin{table}[H]
\TableContent{
    \centering\small
    \setlength{\tabcolsep}{4pt}
    \begin{tabular}{l l cc cc}
  \toprule
  & & \multicolumn{2}{c}{Stopping Game} & \multicolumn{2}{c}{Routing Game} \\
  \cmidrule(lr){3-4} \cmidrule(lr){5-6}
  Architecture & Baseline & $\mathrm{H}$ & $\mathrm{H}_{\mathrm{surv}}$ & $\mathrm{H}$ & $\mathrm{H}_{\mathrm{surv}}$ \\
  \midrule
    \multirow{4}{*}{VGG-16} & clean-intraclass   & $0.513 \pm 0.071$ & $1.000 \pm 0.000$ & $1.000 \pm 0.000$ & $1.000 \pm 0.000$ \\
                               & clean-interclass   & $0.607 \pm 0.032$ & $1.000 \pm 0.000$ & $1.000 \pm 0.000$ & $1.000 \pm 0.000$ \\
                               & clean-vs-random    & $0.663 \pm 0.020$ & $1.000 \pm 0.000$ & $1.000 \pm 0.000$ & $1.000 \pm 0.000$ \\
                               & random-vs-random   & $0.320 \pm 0.019$ & $1.000 \pm 0.000$ & $1.000 \pm 0.000$ & $1.000 \pm 0.000$ \\
    \midrule
    \multirow{4}{*}{ResNet-50} & clean-intraclass   & $0.272 \pm 0.056$ & $0.999 \pm 0.000$ & $0.786 \pm 0.044$ & $1.000 \pm 0.000$ \\
                               & clean-interclass   & $0.268 \pm 0.045$ & $1.000 \pm 0.000$ & $0.902 \pm 0.033$ & $1.000 \pm 0.000$ \\
                               & clean-vs-random    & $0.269 \pm 0.035$ & $1.000 \pm 0.000$ & $0.943 \pm 0.010$ & $1.000 \pm 0.000$ \\
                               & random-vs-random   & $0.152 \pm 0.008$ & $0.992 \pm 0.000$ & $0.739 \pm 0.011$ & $0.999 \pm 0.000$ \\
    \midrule
    \multirow{4}{*}{ViT-B/16} & clean-intraclass   & $0.422 \pm 0.023$ & $0.969 \pm 0.008$ & $0.789 \pm 0.067$ & $1.000 \pm 0.000$ \\
                               & clean-interclass   & $0.435 \pm 0.017$ & $0.974 \pm 0.005$ & $0.882 \pm 0.010$ & $1.000 \pm 0.000$ \\
                               & clean-vs-random    & $0.439 \pm 0.014$ & $0.963 \pm 0.006$ & $0.891 \pm 0.007$ & $1.000 \pm 0.000$ \\
                               & random-vs-random   & $0.377 \pm 0.013$ & $0.867 \pm 0.013$ & $0.476 \pm 0.018$ & $0.997 \pm 0.000$ \\
  \bottomrule
\end{tabular}

    }%
\caption[Reference baselines for the input-noise sweep]{Reference baselines for the input-noise sweep, averaged over $R$ pairs drawn from the
    $565$-image ImageNet-S test split (clean rows: $R=200$ sampled pairs; clean-vs-random: $R=565$
    images paired with a fresh Gaussian each; random-vs-random: $R=200$ Gaussian pairs); entries are
    mean $\pm$ standard deviation over pairs. \textbf{clean-intraclass}: $\mathrm{H}(\Pi_A(x_1),
        \Pi_A(x_2))$ with $x_1, x_2$ from the same ground-truth class; \textbf{clean-interclass}: same,
    with distinct classes; \textbf{clean-vs-random}: one side clean, other Gaussian;
    \textbf{random-vs-random}: both sides Gaussian.}
    \label{ptnb:tab:hellinger-input-random}
\end{table}
\par\medskip
\section{Randomization Sanity Checks}\label{ptnb:app:sanity-check}\label{subsec:hellinger_randomization}

This section supports the cascading-randomization experiments of Section~\ref{ptnb:sec:hellinger-sanity}: \S\ref{ptnb:sec:lmp-randomization} extends the protocol of~\textcite{Adebayo2018sanity} to trajectory Hellinger diagnostics (Section~\ref{ptnb:app:path-divergence}); \S\ref{ptnb:app:sanity-heatmap} tests main-paper attribution maps; \S\ref{ptnb:sec:lmp-rand-double-random} gives the double-random reference; and \S\ref{ptnb:sec:per-image-orthogonality} compares trajectory and map diagnostics per image.

\subsection{Cascading Randomization in Trajectory Space}\label{ptnb:sec:lmp-randomization}

Replacing a weight matrix in the randomization test of~\textcite{Adebayo2018sanity} yields two kernels on the same graph; the Hellinger pass of Section~\ref{ptnb:app:path-divergence} quantifies their trajectory-space change. With its geometric-mean kernel $\widetilde{T}^{(l)}$ and overlap mass $g^{(l)}(s)$, $\mathrm{BC}^{(l)} = \sum_{s \in \mathcal{S}^{(l)}_+} g^{(l)}(s)$; $\mathrm{BC}^{(l)}_{\mathrm{ord}}$ restricts this sum to ordinary states $s \in \mathcal{S}^{(l)}$.

The Stopping-Game contraction bound assumes a conditional alive probability bounded
away from one; the Routing-Game bound assumes a uniform row-overlap bound strictly below one.
The precise hypotheses and proofs are in Theorems~\ref{ptnb:thm:rand-adf}
and~\ref{ptnb:thm:rand-routing}; the cemetery-floor analysis is in
Remark~\ref{ptnb:rem:adf-cemetery-floor}.

The discounted occupation measures at the input layer recover gradient and $\alpha\beta$-\ac{lrp} maps (Theorems~\ref{ptnb:thm:deriv-occ} and~\ref{ptnb:thm:ab-recovery}). The undiscounted endpoint marginal $m^{(0)}_M$ is generally different: Hellinger compares the trajectory routing laws and does not include the path discounts. Under their respective contraction hypotheses, Theorems~\ref{ptnb:thm:rand-adf}--\ref{ptnb:thm:rand-routing} show that $\Pi_M$ changes drastically under randomization, while Table~\ref{tab:sanity-check} shows that the resulting attribution maps can remain similar. \ch{The two
diagnostics therefore detect different aspects of parameter sensitivity: the map-level test
of~\textcite{Adebayo2018sanity} asks whether the explanation delivered to the user depends on the
learned parameters, the trajectory-level test whether its normalized routing does. A map can fail
the first while its routing passes the second; this discrepancy does not isolate a fixed projection
loss, since the discounts can also change across models, and it does not make such a
map adequate for model debugging~\parencite{Adebayo2018sanity,sixt2020whenexplanationslie}.}

\subsection{Heatmap similarity at the cascade endpoint}\label{ptnb:app:sanity-heatmap}

We apply the cascading randomization test of \textcite{Adebayo2018sanity} to Table~\ref{tab:sanity-check}'s attribution methods. Starting pretrained, we cumulatively re-initialize convolutional and linear layers output-to-input: weights Kaiming normal~\parencite{He2015delving}, biases $\mathcal{N}(0, 0.01^2)$. BatchNorm (ResNet-50), LayerNorm, positional embeddings, and class token (ViT-B/16) retain pretrained values. The ViT attribution cascade randomizes decomposed query/key/value projections; the Hellinger cascade (Section~\ref{ptnb:sec:hellinger-cascading}) retains packed query/key/value parameters. At each depth, we compare original and randomized models' attributions for every test image. The test asks whether maps change, without requiring monotonic similarity decay. Persistent similarity can reflect input structure or architecture.

\paragraph{Metrics.}
We report complementary map-similarity measures:
\begin{itemize}
    \item \textbf{Spearman $\rho(|a|)$}: rank correlation of absolute attributions, measuring agreement in pixel importance. Absolute values accommodate differing attribution sign conventions; the correlation itself remains signed.
    \item \textbf{HOG $r$}: Pearson correlation of Histograms of Oriented Gradients (HOG) descriptors~\parencite{Dalal2005hog}, measuring local-gradient-orientation structure.
\end{itemize}
We omit SSIM. \textcite{Binder2022shortcomings} show that input-dependent structure and activation scales can survive top-down randomization, while low similarity can reflect noise; this noise critique also covers Spearman correlation of absolute attributions. Values average over the 565-image ImageNet-S test split (50 classes). Constant maps receive similarity $0$ by convention: at ResNet-50's cascade endpoint, LIME and KernelSHAP outputs no longer depend on superpixels. Table~\ref{tab:sanity-check} reports this endpoint, with all tested convolutional and linear layers re-initialized and the parameter groups above retained.

\paragraph{Discussion.}
Gradient and Grad-CAM stay closest to zero on the \acp{cnn}, as does LayerCAM on VGG-16; on ViT-B/16, Gradient, SmoothGrad and \ac{ptnbGAE} do. These scores record map changes under the tested intervention, not faithfulness. LRP-$\varepsilon$ retains the most VGG-16 rank structure ($\sanityVggLRPepsRho$). ResNet-50's \ac{lrp} rules (LRP-$\gamma$ $\sanityResLRPgammaRho$, LRP-$\varepsilon$ $\sanityResLRPepsRho$) approach Integrated Gradients and LayerCAM, which retain most there. LIME and KernelSHAP retain HOG correlation about $0.86$ with low Spearman correlation on VGG-16 and ViT-B/16: input-derived superpixels preserve edge structure. On ResNet-50, endpoint outputs no longer depend on superpixels, yielding constant maps ($^{\ddagger}$). The \ac{rg} rows lie between Gradient and LRP-$\varepsilon$ on VGG-16, slightly below \ac{lrp} on ResNet-50, and highest in ViT-B/16 Spearman correlation ($\sanityVitRGanchorRho$ anchor, $\sanityVitRGlocRho$ selected mode). Grad-CAM++ on ResNet-50 and \ac{ptnbTIBAV} on ViT-B/16 become anti-correlated, rather than unrelated. Section~\ref{ptnb:sec:lmp-randomization} complements map similarity with trajectory-distribution divergence.

\begin{table}[H]
\centering
\scriptsize
\caption[Parameter randomization sanity check]{Parameter randomization sanity check~\parencite{Adebayo2018sanity}. Each cell reports the average similarity between the attribution of the original pretrained model and the attribution of a fully randomized copy (every convolutional and linear layer re-initialized; normalization layers, the attention's query/key/value projection and the ViT embeddings stay pretrained). Values near zero indicate that the method detects the parameter change; the Spearman correlation is signed, so negative values occur. Bold marks the value closest to zero per metric within each architecture, among maps that are not constant; $^{\ddagger}$: the map is constant at full randomization (similarity $0$ by convention). Rows: the reported configuration of every method (Section~\ref{ptnb:sec:experiments}).}
\label{tab:sanity-check}
\begin{tabular}{l c c}
\toprule
Method & Spear.\ $\rho(|a|)$ & HOG $r$ \\
\midrule
\multicolumn{3}{l}{\textbf{VGG-16}} \\
\midrule
Gradient & 0.013 & \textbf{0.009} \\
SmGrad & 0.084 & 0.101 \\
IntGrad & 0.207 & 0.348 \\
DeepLIFT & 0.216 & 0.340 \\
LRP-$\varepsilon$ & 0.341 & 0.457 \\
LRP-$\gamma$ & 0.092 & 0.370 \\
GCAM & \textbf{-0.007} & 0.069 \\
GCAM++ & 0.180 & 0.129 \\
LCAM & -0.011 & 0.052 \\
LIME & 0.257 & 0.864 \\
KernelSHAP & 0.075 & 0.868 \\
SG ($\lambda{=}{-}1,\sigma^2{=}5$, gate on) & 0.289 & 0.456 \\
RG anchor & 0.299 & 0.535 \\
RG ($\tau{=}0.3,\varepsilon{=}0.1$) & 0.284 & 0.287 \\
\midrule
\multicolumn{3}{l}{\textbf{ResNet-50}} \\
\midrule
Gradient & 0.148 & 0.134 \\
SmGrad & 0.233 & 0.214 \\
IntGrad & 0.277 & 0.424 \\
DeepLIFT & 0.224 & 0.351 \\
LRP-$\varepsilon$ & 0.250 & 0.393 \\
LRP-$\gamma$ & 0.254 & 0.130 \\
GCAM & \textbf{0.020} & \textbf{0.072} \\
GCAM++ & -0.502 & 0.151 \\
LCAM & 0.266 & 0.145 \\
LIME$^{\ddagger}$ & 0.000 & 0.000 \\
KernelSHAP$^{\ddagger}$ & 0.000 & 0.000 \\
RG anchor & 0.217 & 0.344 \\
RG ($\tau{=}0.9,\varepsilon{=}1$) & 0.180 & 0.316 \\
\midrule
\multicolumn{3}{l}{\textbf{ViT-B/16}} \\
\midrule
Gradient & \textbf{-0.003} & 0.041 \\
SmGrad & 0.010 & \textbf{0.038} \\
IntGrad & 0.272 & 0.362 \\
LRP-$\varepsilon$ & 0.165 & 0.300 \\
LRP-$\gamma$ & 0.027 & 0.234 \\
AttnLRP & 0.124 & 0.176 \\
DAVE & 0.178 & 0.323 \\
TIBAV & -0.175 & 0.150 \\
GAE & -0.009 & 0.130 \\
AR & 0.120 & 0.153 \\
LIME & 0.248 & 0.861 \\
KernelSHAP & 0.131 & 0.867 \\
RG anchor & 0.454 & 0.469 \\
RG ($\tau{=}0.4,\lambda_{\mathrm{sm}}{=}{-}10,\sigma^2{=}5,\varepsilon{=}0.1$) & 0.450 & 0.310 \\
\bottomrule
\end{tabular}

\par\smallskip
{\footnotesize Spear.\ $\rho(|a|)$ = Spearman rank correlation of absolute attribution values; HOG $r$ = Pearson correlation of HOG descriptors. RG anchor = the $\alpha\beta$-\ac{lrp} representation ($\tau{=}1$, $\varepsilon{=}0.5$); the second \ac{rg} row and the \ac{sg} row are the localization-selected configurations and list the modes that differ ($\tau$ = temperature; $\lambda$, $\lambda_{\mathrm{sm}}$ = risk aversion; $\sigma^2$ = \ac{ptnbADF} input variance; $\varepsilon$ = \ac{lrp} stabilizer; $\tilde\Phi$ = noisy-observation gate).}
\end{table}


\subsection{Double-Random Empirical Baseline}\label{ptnb:sec:lmp-rand-double-random}

Our cascade-endpoint reference compares two copies of each architecture with independently randomized convolutional and linear layers, retaining the pretrained parameter groups below. This differs from \textcite{Adebayo2018sanity}'s trained-versus-randomized comparison.

\paragraph{Protocol.}
For each architecture, $R = 500$ independent rounds comprise: \emph{(i)} copy the pretrained checkpoint to $A_r, B_r$ and independently re-initialize every \texttt{Conv2d}/\texttt{Linear} layer as at the cascade endpoint. Both retain normalization parameters, packed attention query/key/value projections, positional embedding, and class token. We use checkpoints for all architectures: \texttt{weights=None} gives ViT an identically zero classification head, making trajectory laws trivially indistinguishable. \emph{(ii)} Draw one image with replacement from the $565$-image ImageNet-S test split ($\drDistinctImages$ distinct images among $500$ draws with our fixed seed). \emph{(iii)} Compute plain $\mathrm{H}(\Pi_{A_r}, \Pi_{B_r})$ and survival-conditioned $\mathrm{H}_{\mathrm{surv}}$ for both Stopping and Routing Game kernels, targeting the ground-truth class. Table~\ref{ptnb:tab:double-random-hellinger} reports round-to-round means and standard deviations; variance mixes image selection and random-model-pair realization.

\begin{table}[H]
\TableContent{
    \centering
    \small
    \begin{tabular}{lcccc}
  \toprule
  & \multicolumn{2}{c}{Stopping Game} & \multicolumn{2}{c}{Routing Game} \\
  \cmidrule(lr){2-3} \cmidrule(lr){4-5}
  Architecture & $\mathrm{H}$ & $\mathrm{H}_{\mathrm{surv}}$ & $\mathrm{H}$ & $\mathrm{H}_{\mathrm{surv}}$ \\
  \midrule
    VGG-16 & $0.939 \pm 0.003$ & $1.000 \pm 0.000$ & $1.000 \pm 0.000$ & $1.000 \pm 0.000$ \\
    ResNet-50 & $0.938 \pm 0.003$ & $1.000 \pm 0.000$ & $0.973 \pm 0.002$ & $1.000 \pm 0.000$ \\
    ViT-B/16 & $0.958 \pm 0.002$ & $1.000 \pm 0.000$ & $0.956 \pm 0.001$ & $1.000 \pm 0.000$ \\
  \bottomrule
\end{tabular}

    }%
\caption[Hellinger trajectory distance at the double-random baseline]{Hellinger trajectory distance between two copies of the
    same architecture with independently randomized convolutional and linear layers and the remaining
    pretrained parameters retained, averaged over $R = 500$ rounds (one fresh ImageNet-S test image per round).
    Plain $\mathrm{H}$ concentrates at an architecture-dependent cemetery floor of
    $\drFloorMin$--$\drFloorMax$; conditioned-survival $\mathrm{H}_{\mathrm{surv}}$ saturates at unity.
    The VGG-16 Routing Game cell is the one case in which the plain distance also saturates at
    $1.000$.}
    \label{ptnb:tab:double-random-hellinger}
\end{table}

\paragraph{Interpretation.}
Plain $\mathrm{H}$ concentrates at a model-dependent floor of $\drFloorMin$--$\drFloorMax$,
matching the cemetery-floor prediction of Remark~\ref{ptnb:rem:adf-cemetery-floor}: once both random
kernels produce uncorrelated dead masks, the shared dead-in-both coincidence rate bounds
$\mathrm{BC}^{(0)}$ below by a strictly positive constant and plain $\mathrm{H}$ cannot reach
unity.  The conditioned-survival distance removes the cemetery floor by renormalizing on survival
and reports the two randomized trajectory laws as essentially support-disjoint on
ordinary input states ($\mathrm{H}_{\mathrm{surv}} \approx 1$ throughout).  The
pretrained-vs-randomized Hellinger distance that the cascading sanity check reports should
therefore be read against this architecture-specific empirical reference: values close to the
double-random baseline indicate a trajectory distance comparable to that between two copies
independently randomized over the same layers.

\subsection{Per-Image Association Between Trajectory and Map Diagnostics}\label{ptnb:sec:per-image-orthogonality}

Theorems~\ref{ptnb:thm:rand-adf}--\ref{ptnb:thm:rand-routing} and the double-random baseline of
Table~\ref{ptnb:tab:double-random-hellinger} bound the \emph{aggregate} behavior of the Hellinger
trajectory distance and the pixel-level attribution-map similarity.  A complementary question, motivated by the
argument of \S\ref{ptnb:sec:lmp-randomization}, is whether the two diagnostics agree \emph{per
image}: at fixed randomization depth, does a larger displacement of an image's trajectory
distribution correspond to a more dissimilar attribution map?  Table~\ref{ptnb:tab:per-image-corr} tests this directly.

\paragraph{Protocol.}
For each architecture we pick a cascading step at the same network depth in the Hellinger
and attribution-sanity runs: the layers down to \path{features.26} on VGG-16 and down to
\path{layer4.1.conv3} on ResNet-50 (identical sets of re-initialized layers), and the last
encoder block on ViT-B/16, where the Hellinger cascade re-initializes the MLP and the output
projection of the block and the attribution cascade additionally its $Q$, $K$ and $V$ projections.
The interventions therefore match on the CNNs; on ViT-B/16, the comparison also reflects this
difference in the randomized operators.
The mean $\mathrm{H}$ / $\mathrm{H}_{\mathrm{surv}}$ at these steps can be read off
Figure~\ref{ptnb:fig:hellinger-sanity}.  For each of the $\perImageN$ ImageNet-S test images we then
record four Hellinger scalars ($\mathrm{H}$ plain/surv for the Stopping and Routing Games) and four
attribution-map scalars (Spearman $\rho(|a|)$ and HOG, each computed between the pretrained and
partially-randomized model, for Gradient and for the $\alpha\beta$-\ac{rg} anchor at $\tau{=}1$;
these are the attribution methods whose marginal laws line up with the Hellinger game kernels).
Table~\ref{ptnb:tab:per-image-corr} reports Spearman rank correlations across the $\perImageN$ images.

\begin{table}[H]
\TableContent{
    \centering
    \scriptsize
    \setlength{\tabcolsep}{3.5pt}
    \begin{tabular}{ll cccc cccc}
  \toprule
  & & \multicolumn{4}{c}{Within Hellinger} & \multicolumn{2}{c}{Gradient} & \multicolumn{2}{c}{$\mathrm{RG}$ anchor} \\
  \cmidrule(lr){3-6} \cmidrule(lr){7-8} \cmidrule(lr){9-10}
  Architecture & Hellinger variant
    & $\mathrm{H}_{\mathrm{RG,pl}}$ & $\mathrm{H}_{\mathrm{RG,surv}}$
    & $\mathrm{H}_{\mathrm{SG,pl}}$ & $\mathrm{H}_{\mathrm{SG,surv}}$
    & Spear.\,$\rho(|a|)$ & HOG\,$r$ & Spear.\,$\rho(|a|)$ & HOG\,$r$ \\
  \midrule
\multirow{4}{*}{VGG-16} & $\mathrm{H}_{\mathrm{RG,plain}}$ & $1.00$ & $+1.00$ & $-0.33$ & $+0.85$ & $+0.07$ & $+0.10$ & $-0.01$ & $+0.04$ \\
 & $\mathrm{H}_{\mathrm{RG,surv}}$ & $+1.00$ & $1.00$ & $-0.33$ & $+0.85$ & $+0.07$ & $+0.10$ & $-0.01$ & $+0.04$ \\
 & $\mathrm{H}_{\mathrm{SG,plain}}$ & $-0.33$ & $-0.33$ & $1.00$ & $-0.30$ & $+0.12$ & $+0.11$ & $+0.00$ & $-0.02$ \\
 & $\mathrm{H}_{\mathrm{SG,surv}}$ & $+0.85$ & $+0.85$ & $-0.30$ & $1.00$ & $+0.01$ & $+0.05$ & $-0.10$ & $-0.10$ \\
    \midrule
\multirow{4}{*}{ResNet-50} & $\mathrm{H}_{\mathrm{RG,plain}}$ & $1.00$ & $+0.58$ & $-0.08$ & $+0.70$ & $+0.05$ & $-0.01$ & $-0.08$ & $-0.19$ \\
 & $\mathrm{H}_{\mathrm{RG,surv}}$ & $+0.58$ & $1.00$ & $+0.16$ & $+0.89$ & $+0.15$ & $+0.08$ & $+0.05$ & $-0.05$ \\
 & $\mathrm{H}_{\mathrm{SG,plain}}$ & $-0.08$ & $+0.16$ & $1.00$ & $+0.16$ & $-0.04$ & $+0.07$ & $-0.01$ & $+0.02$ \\
 & $\mathrm{H}_{\mathrm{SG,surv}}$ & $+0.70$ & $+0.89$ & $+0.16$ & $1.00$ & $+0.16$ & $+0.07$ & $-0.01$ & $-0.10$ \\
    \midrule
\multirow{4}{*}{ViT-B/16} & $\mathrm{H}_{\mathrm{RG,plain}}$ & $1.00$ & $+0.96$ & $+0.29$ & $+0.30$ & $+0.02$ & $+0.00$ & $-0.04$ & $-0.08$ \\
 & $\mathrm{H}_{\mathrm{RG,surv}}$ & $+0.96$ & $1.00$ & $+0.25$ & $+0.28$ & $+0.02$ & $+0.02$ & $-0.07$ & $-0.07$ \\
 & $\mathrm{H}_{\mathrm{SG,plain}}$ & $+0.29$ & $+0.25$ & $1.00$ & $+0.99$ & $-0.12$ & $-0.06$ & $-0.03$ & $-0.05$ \\
 & $\mathrm{H}_{\mathrm{SG,surv}}$ & $+0.30$ & $+0.28$ & $+0.99$ & $1.00$ & $-0.12$ & $-0.06$ & $-0.04$ & $-0.05$ \\
  \bottomrule
\end{tabular}

    }%
\caption[Per-image Spearman rank correlations]{Per-image Spearman rank correlations across the $\perImageN$-image ImageNet-S test split
    at the selected cascading steps of the Hellinger and attribution-sanity pipelines. The randomized
    layers match on the CNNs; the ViT-B/16 attribution cascade additionally randomizes $Q$, $K$, and $V$.  Left block:
    within-Hellinger cross-correlations of the plain ($\mathrm{H}_{\cdot,\mathrm{pl}}$) and
    survival-conditioned ($\mathrm{H}_{\cdot,\mathrm{surv}}$) distances under the \ac{rg} and \ac{sg}
    kernels (self-entries shown as $-$; trivially $+1$).  Right block: Hellinger-vs-map
    cross-correlations.  The Hellinger-vs-map block is uniformly $|\rho| \leq \perImageCrossMax$:
    per-image Hellinger and per-image attribution-map similarity have weak observed rank association.}
    \label{ptnb:tab:per-image-corr}
\end{table}

\paragraph{Readings.}
The Hellinger-vs-map block of Table~\ref{ptnb:tab:per-image-corr} is uniformly near zero ($|\rho| \leq \perImageCrossMax$ across all three architectures): trajectory displacement and map drift have weak observed per-image rank association. Statistical independence and nonlinear dependence remain unresolved; ViT-B/16's unequal interventions further limit conclusions about the diagnostics alone.

Within Hellinger, per-image $\mathrm{H}$ gives a stable ordering. When little mass reaches the cemetery, plain and conditioned distances rank images almost identically: $\rho = +0.96$ (\ac{rg}) and $+0.99$ (\ac{sg}) on ViT-B/16. VGG-16 \ac{rg} gives identical values because it never routes to dead-neuron states. For \ac{sg} on \ac{relu} backbones, the normalizations decouple ($\rho = -0.30$ VGG-16, $+0.16$ ResNet-50): image-dependent dead masks change the cemetery share, which renormalization by $1/\sqrt{Z_A Z_B}$ removes. The two games' conditioned distances agree strongly on \ac{relu} backbones ($\rho = +0.85$ VGG-16, $+0.89$ ResNet-50), but weakly on ViT-B/16 ($+0.28$), where the kernels capture partly different per-image structure.

Within an attribution method, Spearman $\rho(|a|)$ and HOG agree well ($\rho = 0.65$--$0.82$, from the same per-image joins), whereas Gradient and $\alpha\beta$-\ac{rg} similarities agree weakly across images ($\rho$ from $-0.10$ to $+0.30$, depending on architecture and metric). Together with the near-zero Hellinger-vs-map block, this supports \S\ref{ptnb:sec:lmp-randomization}: trajectory and pixel diagnostics yield weakly associated per-image rankings here, and the attribution method changes which images a pixel-level sanity check flags.
\par\medskip
\providecommand{\nValGrid}{\textcolor{red}{??}}%

